% Options for packages loaded elsewhere
% Options for packages loaded elsewhere
\PassOptionsToPackage{unicode}{hyperref}
\PassOptionsToPackage{hyphens}{url}
\PassOptionsToPackage{dvipsnames,svgnames,x11names}{xcolor}
%
\documentclass[
  indonesian,
  11pt,
  a4paper,
]{report}
\usepackage{xcolor}
\usepackage[top=2.2cm,bottom=2.2cm,left=2.5cm,right=2.5cm]{geometry}
\usepackage{amsmath,amssymb}
\setcounter{secnumdepth}{5}
\usepackage{iftex}
\ifPDFTeX
  \usepackage[T1]{fontenc}
  \usepackage[utf8]{inputenc}
  \usepackage{textcomp} % provide euro and other symbols
\else % if luatex or xetex
  \usepackage{unicode-math} % this also loads fontspec
  \defaultfontfeatures{Scale=MatchLowercase}
  \defaultfontfeatures[\rmfamily]{Ligatures=TeX,Scale=1}
\fi
\usepackage{lmodern}
\ifPDFTeX\else
  % xetex/luatex font selection
\fi
% Use upquote if available, for straight quotes in verbatim environments
\IfFileExists{upquote.sty}{\usepackage{upquote}}{}
\IfFileExists{microtype.sty}{% use microtype if available
  \usepackage[]{microtype}
  \UseMicrotypeSet[protrusion]{basicmath} % disable protrusion for tt fonts
}{}
\usepackage{setspace}
\makeatletter
\@ifundefined{KOMAClassName}{% if non-KOMA class
  \IfFileExists{parskip.sty}{%
    \usepackage{parskip}
  }{% else
    \setlength{\parindent}{0pt}
    \setlength{\parskip}{6pt plus 2pt minus 1pt}}
}{% if KOMA class
  \KOMAoptions{parskip=half}}
\makeatother
% Make \paragraph and \subparagraph free-standing
\makeatletter
\ifx\paragraph\undefined\else
  \let\oldparagraph\paragraph
  \renewcommand{\paragraph}{
    \@ifstar
      \xxxParagraphStar
      \xxxParagraphNoStar
  }
  \newcommand{\xxxParagraphStar}[1]{\oldparagraph*{#1}\mbox{}}
  \newcommand{\xxxParagraphNoStar}[1]{\oldparagraph{#1}\mbox{}}
\fi
\ifx\subparagraph\undefined\else
  \let\oldsubparagraph\subparagraph
  \renewcommand{\subparagraph}{
    \@ifstar
      \xxxSubParagraphStar
      \xxxSubParagraphNoStar
  }
  \newcommand{\xxxSubParagraphStar}[1]{\oldsubparagraph*{#1}\mbox{}}
  \newcommand{\xxxSubParagraphNoStar}[1]{\oldsubparagraph{#1}\mbox{}}
\fi
\makeatother


\usepackage{longtable,booktabs,array}
\newcounter{none} % for unnumbered tables
\usepackage{calc} % for calculating minipage widths
% Correct order of tables after \paragraph or \subparagraph
\usepackage{etoolbox}
\makeatletter
\patchcmd\longtable{\par}{\if@noskipsec\mbox{}\fi\par}{}{}
\makeatother
% Allow footnotes in longtable head/foot
\IfFileExists{footnotehyper.sty}{\usepackage{footnotehyper}}{\usepackage{footnote}}
\makesavenoteenv{longtable}
\usepackage{graphicx}
\makeatletter
\newsavebox\pandoc@box
\newcommand*\pandocbounded[1]{% scales image to fit in text height/width
  \sbox\pandoc@box{#1}%
  \Gscale@div\@tempa{\textheight}{\dimexpr\ht\pandoc@box+\dp\pandoc@box\relax}%
  \Gscale@div\@tempb{\linewidth}{\wd\pandoc@box}%
  \ifdim\@tempb\p@<\@tempa\p@\let\@tempa\@tempb\fi% select the smaller of both
  \ifdim\@tempa\p@<\p@\scalebox{\@tempa}{\usebox\pandoc@box}%
  \else\usebox{\pandoc@box}%
  \fi%
}
% Set default figure placement to htbp
\def\fps@figure{htbp}
\makeatother



\ifLuaTeX
\usepackage[bidi=basic,shorthands=off,provide=*]{babel}
\else
\usepackage[bidi=default,shorthands=off,provide=*]{babel}
\fi
\ifLuaTeX
  \usepackage{selnolig} % disable illegal ligatures
\fi


\setlength{\emergencystretch}{3em} % prevent overfull lines

\providecommand{\tightlist}{%
  \setlength{\itemsep}{0pt}\setlength{\parskip}{0pt}}



 


\makeatletter
\@ifpackageloaded{caption}{}{\usepackage{caption}}
\AtBeginDocument{%
\ifdefined\contentsname
  \renewcommand*\contentsname{Daftar Isi}
\else
  \newcommand\contentsname{Daftar Isi}
\fi
\ifdefined\listfigurename
  \renewcommand*\listfigurename{Daftar Gambar}
\else
  \newcommand\listfigurename{Daftar Gambar}
\fi
\ifdefined\listtablename
  \renewcommand*\listtablename{Daftar Tabel}
\else
  \newcommand\listtablename{Daftar Tabel}
\fi
\ifdefined\figurename
  \renewcommand*\figurename{Gambar}
\else
  \newcommand\figurename{Gambar}
\fi
\ifdefined\tablename
  \renewcommand*\tablename{Tabel}
\else
  \newcommand\tablename{Tabel}
\fi
}
\@ifpackageloaded{float}{}{\usepackage{float}}
\floatstyle{ruled}
\@ifundefined{c@chapter}{\newfloat{codelisting}{h}{lop}}{\newfloat{codelisting}{h}{lop}[chapter]}
\floatname{codelisting}{Daftar}
\newcommand*\listoflistings{\listof{codelisting}{Daftar Daftar}}
\makeatother
\makeatletter
\makeatother
\makeatletter
\@ifpackageloaded{caption}{}{\usepackage{caption}}
\@ifpackageloaded{subcaption}{}{\usepackage{subcaption}}
\makeatother
\usepackage{bookmark}
\IfFileExists{xurl.sty}{\usepackage{xurl}}{} % add URL line breaks if available
\urlstyle{same}
% fallback for those not using the hyperref driver hyperxmp:
\makeatletter
\@ifundefined{xmpquote}{\newcommand{\xmpquote}[1]{#1}}{}
\makeatother
\hypersetup{
  pdftitle={Executive Summary},
  pdflang={id},
  colorlinks=true,
  linkcolor={blue},
  filecolor={Maroon},
  citecolor={Blue},
  urlcolor={Blue},
  pdfcreator={LaTeX via pandoc}}


\title{Executive Summary}
\usepackage{etoolbox}
\makeatletter
\providecommand{\subtitle}[1]{% add subtitle to \maketitle
  \apptocmd{\@title}{\par {\large #1 \par}}{}{}
}
\makeatother
\subtitle{Zonasi Intervensi Kebijakan UMKM Kota Mataram: Pemetaan
Geospasial untuk Akselerasi Digitalisasi dan Sertifikasi Halal}
\author{}
\date{}
\begin{document}
\maketitle


\setstretch{1.15}
\chapter*{Ringkasan Eksekutif}\label{ringkasan-eksekutif}

Penelitian \emph{Zonasi Intervensi Kebijakan UMKM Kota Mataram: Pemetaan
Geospasial untuk Akselerasi Digitalisasi dan Sertifikasi Halal}
dilaksanakan untuk menyediakan informasi berbasis data yang dapat
mendukung perencanaan pengembangan Usaha Mikro, Kecil, dan Menengah
(UMKM) di Kota Mataram. Penelitian ini berangkat dari kebutuhan untuk
memahami variasi kesiapan digital pelaku UMKM dan keterkaitannya dengan
aktivitas perdagangan elektronik, infrastruktur digital, hambatan
terkait halal, serta kesiapan pemenuhan persyaratan halal.

Penelitian menggunakan pendekatan kuantitatif dengan analisis
deskriptif, pemetaan geospasial, dan pemodelan statistik yang
mempertimbangkan keterkaitan antarwilayah. Data UMKM digunakan untuk
menggambarkan karakteristik dan persebaran usaha, sedangkan informasi
survei digunakan untuk menyusun indeks kesiapan digital serta indikator
yang berkaitan dengan aspek halal. Analisis dilanjutkan dengan pemodelan
\emph{Conditional Autoregressive} (CAR) untuk mengkaji hubungan antara
kesiapan digital dan empat kelompok prediktor, dengan mempertimbangkan
struktur spasial wilayah kelurahan.

Hasil penelitian menghasilkan peta interaktif UMKM Kota Mataram yang
menyajikan informasi lokasi dan karakteristik usaha. Selain itu,
penelitian menghasilkan indeks kesiapan digital (\emph{Digital Readiness
Index} atau DRI), analisis hubungan antara kesiapan digital dan
faktor-faktor yang relevan, serta informasi spasial yang dapat digunakan
sebagai bahan awal untuk mengidentifikasi kebutuhan pendampingan menurut
wilayah.

Hasil pemodelan menunjukkan hubungan positif antara aktivitas
perdagangan elektronik dan indeks kesiapan digital UMKM. Kesiapan halal
juga menunjukkan arah hubungan positif dalam estimasi model. Sementara
itu, estimasi hubungan infrastruktur digital dan hambatan terkait halal
masih mengandung ketidakpastian sehingga perlu ditafsirkan secara
hati-hati. Temuan tersebut menunjukkan pentingnya mempertimbangkan
aktivitas digital dan kesiapan pemenuhan persyaratan halal dalam
perencanaan pengembangan kapasitas UMKM, tanpa mengasumsikan bahwa
seluruh faktor telah terbukti memiliki pengaruh yang pasti.

Kemanfaatan utama penelitian terletak pada penyediaan instrumen
informasi dan dasar analitis bagi pemerintah daerah, lembaga pendamping
UMKM, dan pemangku kepentingan lainnya. Peta interaktif dapat digunakan
untuk mengeksplorasi persebaran usaha dan karakteristik wilayah,
sedangkan hasil analisis dapat menjadi bahan pertimbangan dalam menyusun
prioritas pendampingan, merancang program peningkatan kapasitas digital,
dan mengintegrasikan edukasi halal dengan penguatan kualitas produk.

Diseminasi hasil penelitian turut dilakukan melalui forum diskusi
kelompok terarah (\emph{Focus Group Discussion} atau FGD) pada 13
Oktober 2026 di Mataram. Diskusi melibatkan unsur perguruan tinggi,
pemangku kepentingan terkait UMKM, dan pelaku usaha. Salah satu isu yang
disampaikan oleh pelaku UMKM adalah kebutuhan informasi mengenai proses
pengurusan sertifikasi halal dan penjaminan kualitas produk. Masukan
tersebut memperkuat relevansi agenda pendampingan yang menghubungkan
digitalisasi usaha dengan kesiapan pemenuhan persyaratan halal.

Sebagai tindak lanjut, hasil penelitian dapat dimanfaatkan untuk
mendukung perencanaan program berbasis wilayah, penguatan literasi dan
aktivitas perdagangan elektronik, pendampingan sertifikasi halal, serta
pemantauan perkembangan kesiapan digital UMKM. Pelaksanaannya perlu
disesuaikan dengan kebutuhan aktual di lapangan, disertai verifikasi
data dan koordinasi dengan instansi yang berwenang.

Secara keseluruhan, penelitian ini memberikan kontribusi melalui
pengembangan peta interaktif, indeks kesiapan digital, analisis
statistik dan spasial, serta rekomendasi kebijakan yang dapat menjadi
dasar pendukung pengambilan keputusan berbasis data. Hasil tersebut
diharapkan dapat memperkuat sinergi antara perguruan tinggi, pemerintah
daerah, lembaga pendamping, dan pelaku UMKM dalam mendorong pengembangan
usaha yang lebih adaptif terhadap digitalisasi dan lebih siap memenuhi
persyaratan sertifikasi halal.

\textbf{Luaran utama:} Peta interaktif UMKM Kota Mataram, indeks
kesiapan digital UMKM, hasil pemodelan statistik dan spasial, serta
rekomendasi pemanfaatan untuk perencanaan intervensi kebijakan.

\textbf{Tautan peta interaktif:}
\url{https://indiraputeri.github.io/dudi26/eda\#peta-interaktif-umkm-kota-mataram}

\textbf{Kata kunci:} UMKM, kesiapan digital, perdagangan elektronik,
sertifikasi halal, pemetaan geospasial, model CAR, kebijakan berbasis
data.




\end{document}
